\section*{BLOQUE V. VALORACIÓN DEL DEPARTAMENTO}
\addcontentsline{toc}{section}{BLOQUE V. VALORACIÓN DEL DEPARTAMENTO}

\subsection*{5.1. Otros méritos docentes}
\addcontentsline{toc}{subsection}{5.1. Otros méritos docentes}

\begin{enumerate}[label=\alph*)]

  \item \textbf{Publicaciones de carácter docente y actividades docentes y de divulgación}\addcontentsline{toc}{subsubsection}{a. Publicaciones de carácter docente y actividades docentes y de divulgación}
  
        \begin{itemize}
        
            \item Participación y representación del Instituto Universitario de Investigación en Lingüística Aplicada ILA) en la Noche Europea de los Investigadores celebrada en Cádiz. \hfill 30/09/2017
            
            \item Impartición de la conferencia divulgativa “Diez preguntas sobre el lenguaje” en el marco de la actividad Clases Aplicadas, destinada a alumnado de Secundaria y organizada por el Vicerrectorado de Alumnado – Orientación Universitaria en el IES Santa Isabel de Hungría (Jerez de la Frontera). \hfill 24/01/2023
            
            \item Impartición de la conferencia divulgativa “Diez preguntas sobre el lenguaje” en el marco de la actividad Clases Aplicadas, destinada a alumnado de Secundaria y organizada por el Vicerrectorado de Alumnado – Orientación Universitaria en el IES Elena García Armada (Jerez de la Frontera). \hfill 29/03/2023
            
            \item Participación y representación del Instituto Universitario de Investigación en Lingüística Aplicada ILA) en la Noche Europea de los Investigadores celebrada en Jerez de la Frontera. \hfill 29/09/2023

            \item Diseño y organización de la actividad formativa \textit{"¡Habla ahora o calla para siempre! Trivial de Lingüística"}, celebrada en la Facultad de Filosofía y Letras de la Universidad de Cádiz en el marco de la Semana de las Letras: 2h. \hfill 11/03/2024
            
            \item Participación activa en representación de la Universidad de Cádiz en las XIV Jornadas de Orientación Universitaria, dirigidas a alumnado de Bachillerato y Ciclos Formativos de Grado Superior de la provincia de Cádiz. \hfill 19/03/2024
            
            \item Participación como investigadora en la actividad de divulgación científica «Divulgación científica en la III Feria de Formación Profesional en el Campus de Jerez», organizada por el Vicerrectorado de Investigación y Transferencia de la Universidad de Cádiz. \hfill 23–25/04/2024
            
            \item Participación y representación del Instituto Universitario de Investigación en Lingüística Aplicada ILA) en la Noche Europea de los Investigadores celebrada en Jerez de la Frontera. \hfill 27/09/2024

            \item Participación activa en representación de la Universidad de Cádiz en las XV Jornadas de Orientación Universitaria, dirigidas a alumnado de Bachillerato y Ciclos Formativos de Grado Superior de la provincia de Cádiz. \hfill 9/04/2025
            
            \item Participación como investigadora en la V edición del concurso de divulgación científica \textit{\#HiloTesis: tu tesis doctoral en redes sociales}, organizado por la CRUE en colaboración con la RedDivulga y la Fundación Ignacio Larramendi. \hfill 17/06/2025
            
            \item Participación y representación del Instituto Universitario de Investigación en Lingüística Aplicada ILA) en la Noche Europea de los Investigadores celebrada en San Fernando. \hfill 26/09/2025

            \item Diseño y organización de la actividad formativa \textit{"¡Habla ahora o calla para siempre! Trivial de Lingüística"}, celebrada en la Facultad de Filosofía y Letras de la Universidad de Cádiz en el marco de la Semana de las Letras: 2h. \hfill 11/04/2025

            \item Impartición de tres conferencias divulgativas “El estudio del lenguaje y sus aplicaciones. El grado de Lingüísticas y Lenguas Aplicadas de la UCA”” en el marco de la actividad Clases Aplicadas, destinada a alumnado de Secundaria y organizada por el Vicerrectorado de Alumnado – Orientación Universitaria en el IES Manuel de Falla (Puerto Real). \hfill 4/12/2025

        
        \end{itemize}

  \item \textbf{Talleres, cursos de formación y seminarios especializados}\addcontentsline{toc}{subsubsection}{b. Talleres, cursos de formación y seminarios especializados}

  \begin{itemize}
  \item \textbf{Cursos de doctorado}
  \begin{itemize}[label=\textbullet]
  \item Curso de formaci\'on: \textit{Metodolog\'ia y t\'ecnicas de documentaci\'on en investigaci\'on ling\"u\'istica}\\
  \textbf{Centro:} Escuela de Doctorado de la Universidad de C\'adiz\\
  \textbf{Fecha de finalizaci\'on:} 24/02/2022\\
  \textbf{N.º horas:} 125

  \item Curso de formaci\'on: \textit{B\'usqueda, gesti\'on y comunicaci\'on de la informaci\'on cient\'ifica}\\
  \textbf{Centro:} Escuela de Doctorado de la Universidad de C\'adiz\\
  \textbf{Fecha de finalizaci\'on:} 3/05/2022\\
  \textbf{N.º horas:} 35

  \item Curso de formaci\'on: \textit{Iniciaci\'on a los estudios de Doctorado}\\
  \textbf{Centro:} Escuela de Doctorado de la Universidad de C\'adiz\\
  \textbf{Fecha de finalizaci\'on:} 8/06/2022\\
  \textbf{N.º horas:} 25

  \item Curso de especializaci\'on: \textit{Implicaciones ling\"u\'isticas del contacto intercultural}\\
  \textbf{Centro:} Escuela de Doctorado de la Universidad de C\'adiz\\
  \textbf{Fecha de finalizaci\'on:} 6/06/2023\\
  \textbf{N.º horas:} 125

    \item Curso de especializaci\'on: \textit{Avances en la investigaci\'on del procesamiento ling\"u\'istico y los trastornos del lenguaje}\\
  \textbf{Centro:} Escuela de Doctorado de la Universidad de C\'adiz\\
  \textbf{Fecha de finalizaci\'on:} 30/05/2024\\
  \textbf{N.º horas:} 125

  \item Curso de especializaci\'on: \textit{Directrices actuales de la Ling\"u\'istica en el \'ambito de la comunicaci\'on y las tecnolog\'ias}\\
  \textbf{Centro:} Escuela de Doctorado de la Universidad de C\'adiz\\
  \textbf{Fecha de finalizaci\'on:} 27/06/2024\\
  \textbf{N.º horas:} 125
  \end{itemize}

  
  \item \textbf{Cursos de Formación del PDI}
  \begin{itemize}[label=\textbullet]

  \item \textit{Excel nivel básico. Curso teórico/práctico}\\
  \textbf{Centro:} Universidad de Cádiz\\
  \textbf{Fecha de finalización:} 10/11/2023\\
  \textbf{N.º horas:} 25

  \item \textit{Investigación cuantitativa básica con Excel para CCSS, Educación y Salud orientada a la publicidad}\\
  \textbf{Centro:} Universidad de Cádiz\\
  \textbf{Fecha de finalización:} 01/12/2023\\
  \textbf{N.º horas:} 25

  \item \textit{Estadística aplicada para investigadores con SPSS}\\
  \textbf{Centro:} Universidad de Cádiz\\
  \textbf{Fecha de finalización:} 15/12/2023\\
  \textbf{N.º horas:} 25

  \item \textit{El libro electrónico}\\
  \textbf{Centro:} Universidad de Cádiz\\
  \textbf{Fecha de finalización:} 15/12/2023\\
  \textbf{N.º horas:} 15

  \item \textit{Google App Script: el lenguaje que sirve de pegamento para las herramientas de Google Suite}\\
  \textbf{Centro:} Universidad de Cádiz\\
  \textbf{Fecha de finalización:} 09/02/2024\\
  \textbf{N.º horas:} 25

  \item \textit{Textos científicos con LaTeX en Overleaf}\\
  \textbf{Centro:} Universidad de Cádiz\\
  \textbf{Fecha de finalización:} 09/02/2024\\
  \textbf{N.º horas:} 25

  \item \textit{Excel nivel intermedio. Curso teórico/práctico}\\
  \textbf{Centro:} Universidad de Cádiz\\
  \textbf{Fecha de finalización:} 23/02/2024\\
  \textbf{N.º horas:} 25

  \item \textit{Genera el futuro: aprendiendo inteligencia artificial generativa}\\
  \textbf{Centro:} Universidad de Cádiz\\
  \textbf{Fecha de finalización:} 26/05/2024\\
  \textbf{N.º horas:} 25
  \end{itemize}

\item \textbf{Cursos, jornadas y seminarios}
  \begin{itemize}[label=\textbullet]
    \item Jornadas: \textit{II Jornadas de Divulgación "Las ciencias del
lenguaje al alcance de todos"}\\
  \textbf{Centro:} Universidad de Cádiz – Instituto Universitario de Investigación en Lingüística Aplicada (ILA)
  \hfill 14-15/11/2019\\
  \textbf{N.º horas:} 3
    
    \item Curso: \textit{Plan Integral de Formación para el Empleo (PIFE) - Orientación profesional}\\
  \textbf{Centro:} Universidad de Cádiz
  \hfill 3/11/2020-1/12/2020\\
  \textbf{N.º horas:} 25

    \item Jornadas: \textit{III Jornadas de Divulgación "Las ciencias del
lenguaje al alcance de todos"}\\
  \textbf{Centro:} Universidad de Cádiz – Instituto Universitario de Investigación en Lingüística Aplicada (ILA)
  \hfill 4-5/11/2021\\
  \textbf{N.º horas:} 3
    
    \item Congreso: \textit{2º Congreso Internacional de Jóvenes Investigadores en Lingüística Cognitiva (YRCL2021)}\\
  \textbf{Centro:} Universidad de Alcalá
  \hfill 25–26/11/2021\\
  \textbf{N.º horas:} 20

  \item Seminario: \textit{Metáfora y metonimia, semántica, pragmática y cognición}\\
  \textbf{Centro:} Universidad de Cádiz – Instituto Universitario de Investigación en Lingüística Aplicada (ILA)
  \hfill 14/12/2021\\
  \textbf{N.º horas:} 2.5

  \item Seminario: \textit{El yanito entre la población joven de Gibraltar: bilingüismo, identidad y actitudes lingüísticas}\\
  \textbf{Centro:} Universidad de Cádiz – Instituto Universitario de Investigación en Lingüística Aplicada (ILA)
  \hfill 25/01/2022\\
  \textbf{N.º horas:} 2.5

  \item Jornada: \textit{Recepción para nuevos doctorandos UCA}\\
  \textbf{Centro:} Universidad de Cádiz
  \hfill 01/02/2022\\
  \textbf{N.º horas:} 5

  \item \textit{Jornadas: Doctorales específicas (Lingüística)}\\
  \textbf{Centro:} Universidad de Cádiz
  \hfill 02/02/2022\\
  \textbf{N.º horas:} 6

  \item Seminario de especialización: \textit{Terminología aplicada y documentación terminográfica}\\
  \textbf{Centro:} Universidad de Cádiz – Instituto Universitario de Investigación en Lingüística Aplicada (ILA)
  \hfill 24/02/2022\\
  \textbf{N.º horas:} 4

  \item Curso: \textit{Introducción al diseño de estudios experimentales del lenguaje y su aplicación a la electroencefalografía}\\
  \textbf{Centro:} Universidad de Cádiz – Instituto Universitario de Investigación en Lingüística Aplicada (ILA)
  \hfill 26/04/2022\\
  \textbf{N.º horas:} 10

  \item Jornada: \textit{HEALTHCOM’22. Comunicación, terminología y salud. La evolución del conocimiento en salud, un reto para la comunicación y el lenguaje}\\
  \textbf{Centro:} Universidad Pompeu Fabra
  \hfill 19/05/2022\\
  \textbf{N.º horas:} 9

  \item Curso: \textit{DispoCen: programa para el cálculo de la disponibilidad y centralidad léxica}\\
  \textbf{Centro:} Universidad de Cádiz – Instituto Universitario de Investigación en Lingüística Aplicada (ILA)
  \hfill 17/10/2022\\
  \textbf{N.º horas:} 2

  \item Jornadas: \textit{Lingüística y Cognición}\\
  \textbf{Centro:} Universidad de Cádiz
  \hfill 19–24/10/2022\\
  \textbf{N.º horas:} 20

  \item Taller: \textit{Conceptos estadísticos aplicados a la lingüística de corpus}\\
  \textbf{Centro:} Universidad de León
  \hfill 09/11/2022\\
  \textbf{N.º horas:} 3

  \item Taller: \textit{Herramientas automáticas para la revisión lingüística y la redacción en el lenguaje claro}\\
  \textbf{Centro:} Universidad de Cádiz
  \hfill 22/11/2022\\
  \textbf{N.º horas:} 2

  \item Seminario: \textit{2º Seminario Internacional sobre el Discurso en Twitter}\\
  \textbf{Centro:} Universidad Politécnica de Valencia
  \hfill 30/11–15/12/2022\\
  \textbf{N.º horas:} 4

  \item Jornadas: \textit{Jornadas Doctorales EDUCA y EIDEMAR}\\
  \textbf{Centro:} Universidad de Cádiz
  \hfill 14/12/2022\\
  \textbf{N.º horas:} 5

  \item Jornadas: \textit{Jornadas Doctorales EDUCA y EIDEMAR}\\
  \textbf{Modalidad con póster:} \textit{Terminología y Lingüística computacional: Extracción semi-automática bilingüe de los usos terminológicos de la Lingüística de corpus}\\
  \textbf{Centro:} Universidad de Cádiz
  \hfill 14/12/2022\\
  \textbf{N.º horas:} 10

  \item Coloquio: Coloquio de Doctorandos (Modalidad asistencia) \\
  \textbf{Centro:} Universidad de Cádiz
  \hfill 15/02/2023\\
  \textbf{N.º horas:} 2

  \item Jornada: \textit{Lenguaje y síndromes genéticos raros}\\
  \textbf{Centro:} Universidad de Cádiz
  \hfill 10/03/2023\\
  \textbf{N.º horas:} 7

  \item Seminario: \textit{Introducción al guaraní}\\
  \textbf{Centro:} Universidad de Cádiz – Instituto Universitario de Investigación en Lingüística Aplicada (ILA)
  \hfill 4/05/2023\\
  \textbf{N.º horas:} 2

  \item Seminario: \textit{La trayectoria de las emociones como instrumento de persuasión}\\
  \textbf{Centro:} Universidad de Cádiz – Instituto Universitario de Investigación en Lingüística Aplicada (ILA)
  \hfill 25/05/2023\\
  \textbf{N.º horas:} 2

  \item Jornadas: \textit{XVIII Jornadas de Lingüística Conmemorativas del Trigésimo Aniversario del Grupo Semaínein}\\
  \textbf{Centro:} Universidad de Cádiz – Instituto Universitario de Investigación en Lingüística Aplicada (ILA)
  \hfill 24-25/10/2023 y 8-9/10/2023\\
  \textbf{N.º horas:} 20

  \item Seminario: \textit{3º Seminario Internacional sobre el Discurso en Twitter}\\
  \textbf{Centro:} Universidad Politécnica de Valencia
  \hfill 4/12/2023\\
  \textbf{N.º horas:} 4

  \item Jornadas: \textit{I Jornadas de Inteligencia Artificial en la Lingüística y las Humanidades}\\
  \textbf{Centro:} Universidad de Cádiz
  \hfill 14/05/2024\\
  \textbf{N.º horas:} 4

  \item Seminario: \textit{Lingüística experimental: origen y desarrollos actuales}\\
  \textbf{Centro:} \textit{Online}, organizado por Asociación de Jóvenes Lingüistas
  \hfill 3/06/2024\\
  \textbf{N.º horas:} 2

  \item Seminario: \textit{Análisis del discurso con IAT/ML y LogosLink}\\
  \textbf{Centro:} Universidad de Cádiz – Instituto Universitario de Investigación en Lingüística Aplicada (ILA)
  \hfill 16/05/2024\\
  \textbf{N.º horas:} 2

  \item Seminario: \textit{Metadiscurso en pedagogías multilingües: un enfoque comparativo en Humanidades, Ciencias de la Salud y Ciencias Sociales}\\
  \textbf{Centro:} Universidad de Cádiz – Instituto Universitario de Investigación en Lingüística Aplicada (ILA)
  \hfill 30/05/2024\\
  \textbf{N.º horas:} 2

  \item Coloquio: Coloquio de Doctorandos (Modalidad asistencia) \\
  \textbf{Centro:} Universidad de Cádiz
  \hfill 25/07/2024\\
  \textbf{N.º horas:} 2.5

  \item Coloquio: Coloquio de Doctorandos (Modalidad asistencia) \\
  \textbf{Centro:} Universidad de Cádiz
  \hfill 6/11/2024\\
  \textbf{N.º horas:} 3

  \item Seminario: \textit{The Beggining if Half of the Whole: Writing a good introduction for a research paper}\\
  \textbf{Centro:} Universidad de Cádiz – Instituto Universitario de Investigación en Lingüística Aplicada (ILA)
  \hfill 7/05/2025\\
  \textbf{N.º horas:} 2

  \item Seminario: \textit{Evaluación académica en Humanidades: qué exige ANECA y cómo prepararse}\\
  \textbf{Centro:} Online, organizado por Asociación de Jóvenes Lingüistas
  \hfill 30/06/2025\\
  \textbf{N.º horas:} 2.5

  \item Seminario: \textit{Cuidar la voz: hábitos saludables y ejercicios prácticos}\\
  \textbf{Centro:} Universidad de Cádiz – Instituto Universitario de Investigación en Lingüística Aplicada (ILA)
  \hfill 28/07/2025\\
  \textbf{N.º horas:} 2

    \item Jornada: \textit{XXIV Jornada de la Asociación Española de Terminología: «Terminología y Lexicografía: recursos digitales, plataformas tecnológicas y diccionarios especializados}\\
  \textbf{Centro:} Universidad de Cádiz – Facultad de Filosofía y Letras
  \hfill 6/11/2025\\
  \textbf{N.º horas:} 10

  \item Seminario: \textit{I Seminario de HHDD
aplicadas a la investigación en Literatura y Lengua
Españolas}\\
  \textbf{Centro:} Universidad de Málaga – Departamento de Filología Española, Italiana, Románica y
Teoría de la Literatura y Literatura Comparada y Ciencias y
Técnicas Historiográficas
  \hfill 16/12/2025\\
  \textbf{N.º horas:} 2
  
  \end{itemize}

\item \textbf{Cursos y tutoriales reconocidos por organismos oficiales}

\begin{itemize}[label=\textbullet]

  \item Curso: \textit{Corrección profesional}\\
  \textbf{Centro:} Cálamo y Cran -- Universidad Europea de Madrid\\
  \hfill \textbf{Fecha de finalización:} 18/05/2021\\
  \textbf{N.º horas:} 80

  \item Curso: \textit{Shakespeare's Language: Revealing Meanings and Exploring Myths}\\
  \textbf{Centro:} Universidad de Lancaster (Reino Unido)\\
  \hfill \textbf{Fecha de finalización:} 01/07/2021\\
  \textbf{N.º horas:} 12

  \item Curso: \textit{Software Omniscan 12 y operación del escáner cenital OS 12002 Advanced Plus}\\
  \textbf{Centro:} Universidad de Cádiz -- Instituto Universitario de Investigación en Lingüística Aplicada (ILA)\\
  \hfill \textbf{Fecha de finalización:} 30/11/2021\\
  \textbf{N.º horas:} 5

  
    \item Curso: \textit{Natural Language Processing with Classification and Vector Spaces}\\
      \textbf{Centro:} DeepLearning.AI -- Coursera\\
    \textbf{Fecha de finalización:} 26/05/2025 \\
    \textbf{N.º horas:} 30
    \end{itemize}
  \end{itemize}


  \item \textbf{Participación en congresos y jornadas}\addcontentsline{toc}{subsubsection}{c. Participación en congresos y jornadas}
  \begin{itemize}
  
    
\begin{enumerate}[label={[C\arabic*]}]
\item Moyano Moreno, I. (2022). "Detección de relaciones semánticas para la elaboración de fichas terminológicas y metalingüísticas". XIII International Corpus Linguistics Conference (Bérgamo, Italia).

\item Moyano Moreno, I. (2022). "Humanidades digitales y estilometría: Lectura distante, lingüística digital y visualización macroscópica de textos". XXXVI Congreso Internacional Asociación de Jóvenes Lingüistas (AJL) (Sevilla, España).

\item Moyano Moreno, I. (2022). "Humanidades digitales: Revisión teórica y propuesta de ficha  terminológica del término estilo como uso técnico vinculado a la  Estilometría". VI Congreso de Lingüística y Literatura (Santander, España).

\item Moyano Moreno, I. (2022). "Terminología y Lingüística computacional: extracción semi-automática bilingüe de los usos terminológicos de la Lingüística de corpus". Jornadas Doctorales EDUCA y EIDEMAR. Modalidad con Póster (Cádiz, España).

\item Moyano Moreno, I. y García Cabello de Alba, A. (2023). "Extracción semi-automática de términos vinculada a los contenidos y perspectivas actuales de la Fraseología". XXIII Congreso Internacional de la Asociación Alemana de Hispanistas (Graz, Austria). <\href{https://unipub.uni-graz.at/obvugrveroeff/download/pdf/8389514}{Libro de resúmenes, pp. 462-463}>

\item Moyano Moreno, I. y García Cabello de Alba, A. (2023). "Procesamiento y producción flexiva de neologismos en Twitter". 40º Congreso Internacional de la Sociedad Española de Lingüística Aplicada (Mérida, España).

\item Moyano Moreno, I. y García Cabello de Alba, A. (2023). "Propuesta metodológica desde el procesamiento del lenguaje natural para la creación de un corpus de neologismos: El habla coloquial en Twitter". 40º Congreso Internacional de la Sociedad Española de Lingüística Aplicada (Mérida, España).

\item Moyano Moreno, I. (2024). "Extracción semi-automática de neología en Twitter/X: Hacia una base de datos de neologismos potenciales basada en el léxico de la comunidad hispanohablante de streaming". Jornadas Doctorales EDUCA y EIDEMAR. Modalidad con Póster (Cádiz, España).

\item Moyano Moreno, I. (2023). "Extracción semi-automática de términos vinculada a los contenidos actuales de la Lingüística de corpus". XV Congreso de Lingüística General (Madrid, España).

\item Moyano Moreno, I. y Crespo Miguel, M. (2024). "Tendencias actuales en la intersección terminografía-inteligencia artificial". 41º Congreso Internacional de la Sociedad Española de Lingüística Aplicada (Valencia, España).

\item Moyano Moreno, I. (2024). "Hybrid approaches to Automatic Term Extraction in the Linguistics domain". 41º Congreso Internacional de la Sociedad Española de Lingüística Aplicada (Valencia, España).

\item Moyano Moreno, I. (2024). "De dónde venimos y hacia dónde vamos: cuestiones historiográficas, teóricas y terminológicas en torno a la lingüística de corpus". XXXVIII Congreso Internacional de la Asociación de Jóvenes Lingüistas (Santiago de Compostela, España). <\href{https://ilg.usc.gal/ajlsantiago/wp-content/uploads/2025/02/libro_resumos-3.pdf}{Libro de resúmenes, p. 163}>

\item Moyano Moreno, I. y García Cabello de Alba, A. (2024). "Métodos para la recogida de datos en las ciencias del lenguaje: Del corpus a la experimentación". III Congreso Internacional Tránsitos (Cádiz, España).

\item Moyano Moreno, I. (2024). "Detección semi-automática de neología en Twitter/X: Una base de datos de neologismos potenciales basada en el léxico de la comunidad hispanohablante de streaming". V Congreso Internacional de Formación Permanente: NODOS DEL CONOCIMIENTO. <\href{https://cilg2025.ua.es/uploads/site/files/Libro_Resumenes_CILG25.pdf}{Libro de resúmenes, pp. 1808-1809}>

\item Moyano Moreno, I. y García Cabello de Alba, A. (2024). "Del corpus a la experimentación: Métodos para la recogida de datos en las ciencias del lenguaje". V Congreso Internacional de Formación Permanente: NODOS DEL CONOCIMIENTO. <\href{https://cilg2025.ua.es/uploads/site/files/Libro_Resumenes_CILG25.pdf}{Libro de resúmenes, pp. 1820-1821}>

\item Moyano Moreno, I. "Consideraciones sobre la lingüística de corpus y sus implicaciones terminológicas". V Congreso Internacional de Formación Permanente: NODOS DEL CONOCIMIENTO. <\href{https://cilg2025.ua.es/uploads/site/files/Libro_Resumenes_CILG25.pdf}{Libro de resúmenes, pp. 1823-1824}>

\item Moyano Moreno, I. (2024). "Aportaciones del procesamiento el lenguaje natural a la terminología: modelo teórico-computacional para la extracción de términos lingüísticos". Jornadas Doctorales EDUCA y EIDEMAR. Modalidad oral (Cádiz, España).

\item Moyano Moreno, I. (2024). "Interfaces entre lingüística forense y lingüística computacional: la atribución de autoría". II Sesión del Seminario de Jóvenes Investigadores organizado por la asociación \textit{Ubi Sunt?} (Cádiz, España).

\item Crespo Miguel, M., Moyano Moreno, I. y Bernal Ortiz, F. (2025). "Análisis de autoría en textos robóticos: enfoque lingüístico-forense mediante técnicas cualitativas y cuantitativa". XVI Congreso Internacional de Lingüística general (Alicante, España). <\href{https://cilg2025.ua.es/uploads/site/files/Libro_Resumenes_CILG25.pdf}{Libro de resúmenes, p. 148}>

\item García Cabello de Alba, A. y Moyano Moreno, I. (2025). "Miradas que traducen: un estudio con eye-tracking sobre anglicismos y sus equivalentes traductológicos". XXXIX Congreso Internacional de la Asociación de Jóvenes Lingüistas (Vitoria, España). <\href{https://www.ehu.eus/documents/d/heis/-2409-final-ajl-booklet}{Libro de resúmenes, p. 3}>

\item Moyano Moreno, I. (2025). "¿Qué es un término? Un estudio sobre el (des)acuerdo en la anotación de corpus técnicos". XXXIX Congreso Internacional de la Asociación de Jóvenes Lingüistas (Vitoria, España). <\href{https://www.ehu.eus/documents/d/heis/-2409-final-ajl-booklet}{Libro de resúmenes, p. 177}>


\item Moyano Moreno, I. y Bernal Ortiz, F. (2025). "Terminología y modelos generativos del lenguaje: evaluación en tareas de extracción automática". XXIV Jornada AETER: Terminología y Lexicografía: recursos digitales, plataformas tecnológicas y diccionarios especializados (Cádiz, España).

\end{enumerate}
  
\end{enumerate}

\subsection*{5.2. Pertenencia a grupos de investigación}
\addcontentsline{toc}{subsection}{5.2. Pertenencia a grupos de investigación}


\begin{itemize}

      \item Miembro del grupo de investigación HUM-147 "Semaínein" (PAIDI), grupo interuniversitario y consolidado del Plan Andaluz de Investigación subvencionado por la Junta de Andalucía desde 1994. Dirigido por el Prof. Dr. Miguel Casas Gómez (Universidad de Cádiz). \hfill 2023 -- actualidad
      
      \item Miembro del Instituto Universitario de Investigación en Lingüística Aplicada (ILA) de la Universidad de Cádiz (UCA). \hfill 2025 -- actualidad
      
      \item Otras sociedades y asociaciones:
      \begin{itemize}[label=\textbullet]
            \item Miembro y colaboradora activa del laboratorio de Lingüística computacional y digital del Instituto Universitario de Investigación en Lingüística Aplicada (ILA) de la Universidad de Cádiz. \hfill 2022 - actualidad
            \item Miembro de la Asociación Española de Lingüística de Corpus (AELINCO). \hfill 2022 -- actualidad
            \item Miembro de la Sociedad Española de Procesamiento del Lenguaje Natural (SEPLN). \hfill 2025 -- actualidad
        \end{itemize}
\end{itemize}

\subsection*{5.3. Pertenencia a proyectos de Investigación e Innovación Docente}
\addcontentsline{toc}{subsection}{5.3. Pertenencia a proyectos de Investigación e Innovación Docente}



\begin{itemize}
  \item Miembro del equipo de trabajo en el proyecto \textit{Aplicaciones de la lingüística digital al ámbito de la terminología: la creación de  un léxico relacional bilingüe de los usos terminológicos de la semántica léxica  (TerLexWeb)}\\
  \textbf{Entidad financiadora:} Ministerio de Ciencia e Innovación, Agencia Estatal de Investigación\\
  \textbf{Tipo de convocatoria:} Nacional\\
  \textbf{Fecha de inicio:} 01/09/2023 \hfill \textbf{Fecha de finalización:} 31/08/2027\\
  \textbf{Investigador principal:} Miguel Casas Gómez\\

  \item Miembro del equipo de trabajo en el proyecto \textit{Reconocimiento del origen robótico de textos. Automatización de tareas y conocimiento lingüístico (ROBOT-TALK)}\\
  \textbf{Entidad financiadora:} Ministerio de Ciencia e Innovación, Agencia Estatal de Investigación\\
  \textbf{Tipo de convocatoria:} Nacional\\
  \textbf{Fecha de inicio:} 01/09/2023 \hfill \textbf{Fecha de finalización:} 31/08/2026\\
  \textbf{Investigadoras principales:} Ana María Fernández-Pampillón Cesteros y Marianela Fernández Trinidad\\

  \item Miembro del equipo de trabajo en el proyecto \textit{Longitudinal and cross-sectional neurocognitive study of acoustic-perceptual abilities and linguistic function in preterm infants (PretermLA)}\\
  \textbf{Entidad financiadora:} Vicerrectorado de Política Científica y Tecnológica (Universidad de Cádiz)\\
  \textbf{Tipo de convocatoria:} Nacional\\
  \textbf{Fecha de inicio:} 01/01/2024 \hfill \textbf{Fecha de finalización:} 31/12/2025\\
  \textbf{Investigadores principales:} Carmen Varo Varo y Ricardo Hernández Molina\\

\end{itemize}

\subsection*{5.4. Otros méritos de investigación}
\addcontentsline{toc}{subsection}{5.4. Otros méritos de investigación}

\begin{enumerate}[label=\alph*)]
  \item \textbf{Estancias científicas}\addcontentsline{toc}{subsubsection}{a. Estancias científicas}

\begin{itemize}
    \item Estancia científica internacional en el Linguistics Research Center de la Universidad NOVA de Lisboa (CLUNL) \hfill marzo-mayo 2026
    \item Estancia científica nacional en el Departamento de Filología de la Universidad Compluetense de Madrid (UCM) \hfill junio-julio 2026
\end{itemize}


      \item \textbf{Premios y becas de investigación}\addcontentsline{toc}{subsubsection}{b. Premios y becas de investigación}
    \begin{itemize}
              \item Dotación Económica de 2000€ para Personal Investigador Predoctoral 2025, financiada por el Banco Santander y concedida en régimen de concurrencia competitiva, para la financiación de actividades formativas y estancias de investigación. \hfill 2025
  \end{itemize}

    \item \textbf{Evaluación de artículos científicos}\addcontentsline{toc}{subsubsection}{c. Evaluación de artículos científicos}

    \begin{itemize}
    \item Evaluación por pares ciegos en el n. 1 de la revista {Ramas. Revista de lengua española}}, número 1, 2025. {Índice de evaluadores: <\url{https://journals.uco.es/ramas/article/view/18884/16612}>
    \item Evaluación por pares ciegos en el panel "Lingüística computacional, lingüística de corpus y procesamiento del lenguaje natural" en el marco del \textit{43º Congreso Internacional de AESLA 2026} (15-17 de abril de 2026, Universidad de Granada)
    \item Evaluación por pares ciegos en el panel "Lenguas para fines académicos y profesionales" en el marco del \textit{43º Congreso Internacional de AESLA 2026} (15-17 de abril de 2026, Universidad de Granada)
    \end{itemize}

  \item \textbf{Organización de congresos y jornadas}\addcontentsline{toc}{subsubsection}{d. Organización de congresos y jornadas}
\begin{itemize}
  \item Miembro del comité organizador de las \textit{III Jornadas de Divulgación Científica “Las ciencias del lenguaje al alcance de todos”}, Universidad de Cádiz. \hfill 4–5/11/2021

  \item Miembro del comité organizador de las \textit{I Jornadas de Inteligencia Artificial en la Lingüística y las Humanidades}, Universidad de Cádiz. \hfill 14–15/05/2024

  \item Miembro del comité científico de las \textit{I Jornadas de Inteligencia Artificial en la Lingüística y las Humanidades}, Universidad de Cádiz. \hfill 14–15/05/2024

    \item Moderación de mesa de comunicaciones en el \textit{III Congreso Internacional Tránsitos: Identidades y culturas en movimiento. «Humanidades y progreso desde la era analógica a la digital»}, Universidad de Cádiz. \hfill 20–22/11/2024

  \item Miembro del comité organizador de las \textit{II Jornadas de Inteligencia Artificial en la Lingüística y las Humanidades}, Universidad de Cádiz. \hfill 1 y 8/04/2025

  \item Miembro del comité científico de las \textit{II Jornadas de Inteligencia Artificial en la Lingüística y las Humanidades}, Universidad de Cádiz. \hfill 1 y 8/04/2025

    \item Colaboración con el comité organizador de las \textit{XIX Jornadas de Lingüística}, Universidad de Cádiz. \hfill 14 y 15/10/2025

    \item Miembro del comité organizador de las \textit{XXIV Jornada AETER. Terminología y Lexicografía: recursos digitales, plataformas tecnológicas y diccionarios especializados}, Universidad de Cádiz. \hfill 6/11/2025
\end{itemize}

\subsection*{5.5. Prácticas curriculares}
\addcontentsline{toc}{subsection}{5.5. Prácticas curriculares}


\begin{itemize}
\item Realización de prácticas curriculares en la Universidad de Cádiz, en el marco de la asignatura «Prácticas de empresa» del Grado en Lingüística y Lenguas Aplicadas.\\
Nº. horas: 150\\
Calificación: 9 (Sobresaliente) · Créditos superados: 6 ECTS.\hfill 04/11/2019 – 10/02/2020\\
\end{itemize}


\subsection*{5.6. Alumno/a Colaborador/a}
\addcontentsline{toc}{subsection}{5.6. Alumno/a Colaborador/a}

    \begin{itemize}
      \item Curso 2016-2017: Alumna colaboradora en el Área de Lingüística General de la Universidad de Cádiz, en el marco de las investigaciones del proyecto PRESEEA.
      \item Curso 2017-2018: Alumna colaboradora en el Área de Lingüística General de la Universidad de Cádiz, en el marco de las investigaciones del proyecto PRESEEA.
\end{itemize}

\subsection*{5.7. Méritos adicionales}
\addcontentsline{toc}{subsection}{5.7. Méritos adicionales}

\begin{enumerate}[label=\alph*)]
\item \textbf{Becas disfrutadas}
\begin{itemize}
    \item Concesión de una beca Erasmus+ Estudios KA103 en el programa "Modern Languages" de Newcastle University, Newcastle Upon Tyne (Reino Unido) \hfill Curso 2018/2019
    \item Concesión, en régimen de concurrencia competitiva, de una beca predoctoral FPU (Formación del Profesorado Universitario, convocatoria 2021), en el marco del Programa Estatal para Desarrollar, Atraer y Retener Talento del Plan Estatal de Investigación Científica, Técnica y de Innovación 2021-2023. Convocatoria gestionada por la Secretaría General de Universidades (Ministerio de Universidades, Gobierno de España). \hfill 2022
    \item Concesión de una beca para Personal Investigador Predoctoral 2025, financiada por el Banco Santander y concedida en régimen de concurrencia competitiva, para la financiación de actividades formativas y estancias de investigación. \hfill 2025
\end{itemize}

\item \textbf{Otras ayudas disfrutadas}
\begin{itemize}
    \item Ayuda concedida, en régimen de concurrencia competitiva, por la Asociación de Jóvenes Lingüistas (AJL), modalidad A (ayuda completa), para la participación en el "XXXVI Congreso Internacional de la Asociación de Jóvenes Lingüistas", Sevilla (España). \hfill 2022
    \item Ayuda del Plan Propio UCA 2022-2023 de movilidad para la participación en el "XIII Congreso Internacional de Lingüística de Corpus Discurso de especialidad y estudios de traducción desde la óptica del corpus", Bérgamo (Italia). \hfill 2022
    \item Ayuda del Plan Propio UCA 2022-2023 de movilidad para la participación en el "XXIII. Deutscher Hispanistentag", Graz (Austria). \hfill 2023
    \item Ayuda del Contrato Programa 2025 del Departamento de Filología de la Universidad de Cádiz (UCA) para la participación en el "40º Congreso Internacional de la Sociedad Española de Lingüística Aplicada (AESLA)", Mérida (España). \hfill 2023
    \item Ayuda del Plan Propio UCA 2022-2023 para la realización y el aprovechamiento del curso "Implicaciones lingüísticas del contacto intercultural" en el marco del programa de Doctorado en Lingüística. \hfill 2023
    \item Ayuda del Plan Propio UCA 2022-2023 de movilidad para la participación en el "XV Congreso Internacional de Lingüística General", Madrid (España). \hfill 2023
    \item Ayuda del Plan Propio UCA 2022-2023 de movilidad para la participación en el "41 Congreso Internacional de la Asociación Española de Lingüística Aplicada (AESLA)", Valencia (España). \hfill 2024
    \item Ayuda del Plan Propio UCA 2022-2023 para la realización y el aprovechamiento del curso “Avances en la investigación del procesamiento lingüístico y los trastornos del lenguaje” en el marco del programa de Doctorado en Lingüística. \hfill 2024
    \item Ayuda del Plan Propio UCA 2022-2023 para la realización y el aprovechamiento del curso "Directrices actuales de la Lingüística en el ámbito de la comunicación y las tecnologías" en el marco del programa de Doctorado en Lingüística. \hfill 2024
    \item Ayuda del Contrato Programa 2025 del Departamento de Filología de la Universidad de Cádiz (UCA) para la participación en el "XXXIX Congreso Internacional de la Asociación de Jóvenes Lingüistas", Vitoria (España). \hfill 2025
    \item Ayuda Económica para Personal Investigador Predoctoral 2025, financiada por el Banco Santander y concedida en régimen de concurrencia competitiva, para la financiación de actividades formativas y estancias de investigación. \hfill 2025


\end{itemize}

\item \textbf{Conocimiento de lenguas extranjeras}

\begin{table}[h]
\centering
\begin{tabular}{lll}
\textbf{Idioma} & \textbf{Certificación (MCER)} & \textbf{Expedido por} \\
\hline
Inglés  & C1 & University of Cambridge \\
Francés & B1 & Escuela Oficial de Idiomas \\
Alemán  & A2 & Universidad de Cádiz \\

\end{tabular}
\end{table}
\end{enumerate}


